\section{Pair decoding and soundness}
\label{dec:section}

Set
\begin{equation}
 \sigma=\frac1{h^2\ell^2},
 \label{src:decoding-soundness}
\end{equation}
and assume throughout this section that \(\val(\Phi)\leq\sigma\).
The source promise bounds every labeling of \(U\cup V\), so the
decoders below may use arbitrary fixed instance-dependent data.
Their essential restriction is that the label returned at a variable
does not also depend on the sampled constraint occurrence.

Fix a set \(J\subseteq[d]\) of size \(h\) and frozen data
\(\mathbf a\) as in \Cref{res:private-variance}. The remaining
seeds \(c_e\), \(e\in\binom J2\), are independent uniform
elements of \([M]\); the weight blocks \(s_j\), \(j\in J\),
are independent uniforms on \([-1,1]^r\), independent of the
seeds. Let \(u_j\) be position \(j\)'s tuple of questions from
its incident hidden seeds. Together with \(\mathbf a\), this
tuple determines its query \(i_j\).

Consider any fixed total measurable maps
\begin{equation}
 L_j=\mathcal L_j(u_j,s_j)\subseteq\Lambda_{i_j},
 \qquad |L_j|\leq\ell.
 \label{dec:local_lists}
\end{equation}
Each map receives only the indicated own questions and weights;
on an unsupported own-question tuple it returns the empty list.
Order each list by its slots. The lists in \Cref{ana:lists} have
exactly these properties.

For distinct \(j,k\in J\), let \(E_{jk}\) be the event that
some \(I\in\mathcal I\) hits both lists: it contains
\((i_j,\lambda_j)\) and \((i_k,\lambda_k)\) for labels
\(\lambda_j\in L_j\), \(\lambda_k\in L_k\). Define
\begin{equation}
 Z=\max_{I\in\mathcal I}
       \#\{j\in J:\text{ some }\lambda\in L_j
                              \text{ has }(i_j,\lambda)\in I\}.
 \label{dec:hits}
\end{equation}

\begin{lemma}[Decoding a pair]
\label{dec:pair}
For every fixed \(J,\mathbf a\) and fixed list maps satisfying
\Cref{dec:local_lists},
\begin{equation}
 \Pr(E_{jk})\leq\ell^2\sigma
 \quad(j,k\in J,\ j<k),
 \qquad
 \E Z\leq1+\binom h2\ell^2\sigma
       =\frac32-\frac1{2h}<\frac32.
 \label{dec:pair_bound}
\end{equation}
\end{lemma}
\begin{proof}
Fix \(j<k\) and write \(e=\{j,k\}\). Fix all batch weights
and all hidden seeds except \(c_e\). By independence, the
remaining occurrence \(c_e\) is still uniform on \([M]\).
Position \(j\)'s only varying question is \(x_{c_e}\in U\),
and position \(k\)'s only varying question is \(y_{c_e}\in V\).
Every other argument of each list map is now a constant, while the
maps themselves remain fixed. For the maps constructed in
\Cref{ana:lists}, their conditional averages are not recomputed:
the actual values just fixed are not supplied as new arguments to
those averages.

For \(x\in U\), form the tuple \(u_j^{x}\) by putting \(x\)
in the distinguished \(e\)-coordinate of \(u_j\) and retaining
all its other, now fixed, coordinates. Define \(u_k^{y}\) for
\(y\in V\) in the same way. Choose default symbols
\(a_U\in\Sigma_U\) and \(a_V\in\Sigma_V\). For each pair
of ranks \(\rho,\tau\in[\ell]\), define a labeling
\(A_{\rho,\tau}\) of the entire source instance as follows.
At a variable \(x\in U\), evaluate
\(\mathcal L_j(u_j^{x},s_j)\). If rank \(\rho\) is present,
give \(x\) its value in that label; otherwise give it \(a_U\).
At a variable \(y\in V\), use its value in the label at rank
\(\tau\) of \(\mathcal L_k(u_k^{y},s_k)\), or \(a_V\) if
the rank is absent.

This gives one symbol at every source variable. If a substituted
tuple has probability zero, its empty list invokes the same
fallback. If the variable already occurs in another coordinate of
the query, its value in a present label is still unique: a local
label is an assignment on the set of distinct variables in the
scope. Thus repeated variables, empty valid-label sets, and missing
ranks cause no ambiguity. The assigned symbol on \(U\) depends
only on \(x\), and the assigned symbol on \(V\) only on \(y\),
besides fixed data. It does not use the occurrence index or its
projection map. The source promise consequently gives
\begin{equation}
 \Pr_{c\text{ uniform in }[M]}
 \bigl[\pi_c(A_{\rho,\tau}(x_c))
                    =A_{\rho,\tau}(y_c)\bigr]\leq\sigma.
 \label{dec:rank_soundness}
\end{equation}

For every actual remaining occurrence \(c_e\), the reconstructed
tuples \(u_j^{x_{c_e}},u_k^{y_{c_e}}\) are supported under the
original hidden-seed law: the full seed tuple comprising the fixed
other seeds and this occurrence has positive probability in that
product law. They are exactly the own tuples at which the actual
lists are evaluated.

Suppose \(E_{jk}\) occurs and choose the ranks \(\rho,\tau\)
of labels witnessing it. The endpoint \(x_{c_e}\) belongs to
\(P(i_j)\), and \(y_{c_e}\) to \(P(i_k)\). Compatibility
requires the induced labeling to satisfy every source constraint
with both endpoints in the union of the selected scopes. In
particular it satisfies the sampled occurrence \(c_e\). Its
endpoint labels are exactly those used by
\(A_{\rho,\tau}\), so
\[
 \pi_{c_e}(A_{\rho,\tau}(x_{c_e}))
                 =A_{\rho,\tau}(y_{c_e}).
\]
This uses the sampled occurrence even when other occurrences have
the same endpoints. Union bounding \Cref{dec:rank_soundness} over
the \(\ell^2\) rank pairs proves the pair estimate for the fixed
weights and other seeds. It is uniform in those fixed values, so
averaging proves \(\Pr(E_{jk})\leq\ell^2\sigma\). No
conditioning on compatible labels or successful ranks is used.

For every integer \(z\geq0\),
\(z\leq1+\binom z2\), since the difference is
\((z-1)(z-2)/2\geq0\). Choose a selection attaining the maximum
in \Cref{dec:hits}. Each of the pairs among its \(Z\) hit
positions witnesses the corresponding event \(E_{jk}\).
Consequently, pointwise,
\begin{equation}
 Z\leq1+\binom Z2
   \leq1+\sum_{\{j,k\}\in\binom J2}\ind_{E_{jk}}.
 \label{dec:pair_count}
\end{equation}
Taking expectations, using the pair estimate, and substituting
\(\sigma=1/(h^2\ell^2)\) gives \Cref{dec:pair_bound}.
\end{proof}

\begin{corollary}[Soundness]
\label{dec:soundness}
For the graph constructed in \Cref{graph:section}, assume
\eqref{src:restriction-budget}.
If \(\val(\Phi)\leq\sigma\), every independent set has size
strictly less than \(n/m\).
Consequently \(\tau(G)>(1-1/m)n\).
\end{corollary}
\begin{proof}
Suppose an independent set \(\mathcal A\) has size at least
\(n/m=2b_0n\). By the low-norm clique observation at the start of
\Cref{ana:section}, at most one vertex of \(\mathcal A\) has
\(\norm{S_v}\leq t\). Since \(n\geq p>1/b_0\), the number
of its high-norm vertices is at least
\[
 2b_0n-1>b_0n.
\]

\Cref{ana:function} therefore supplies the single fixed odd
Lipschitz function \(F\) and its variance lower bound. Applying
\Cref{res:private-variance} and then the list construction of
\Cref{ana:lists,ana:list-hits} gives fixed \(J,\mathbf a\)
and local lists with
\[
 \E Z\geq\frac{\beta h}{2}=2.
\]
This contradicts \Cref{dec:pair}, which gives \(\E Z<3/2\)
under the same remaining product law.

The lower bound just used averages over the hidden seeds and
continuous batch weights for those fixed \(J,\mathbf a\).
The upper bound fixes further weights and seeds only to decode a
single occurrence and is then averaged back. It does not require
an energy lower bound, or a derivative identity, at every such
further fixing.

We have ruled out every independent set of size at least \(n/m\).
The complement of any vertex cover is independent, so every cover
has size strictly greater than \(n-n/m\), as claimed.
\end{proof}
